% !TEX root = ../../proposal.tex

\section{Phase 3: Representation Of Co-morbid Disease Conditions Using Composable Diffusion Models}
\label{sec:phase_3}
This phase is the clinical extension of the investigation developed in Phases~1 and~2.
Phase~1 motivates why compositional generation can fail when the training distribution induces conflicting concept interactions, while Phase~2 will test whether explicit factorisation can recover held-out composition in a controlled synthetic setting.
Building on that logic, this phase asks whether the same representation-centred account can be applied in a realistic clinical domain.

We study posteroanterior (PA) chest X-rays containing multiple disease conditions and treat co-morbid synthesis as an out-of-distribution composition problem.
If H2 is supported, then this phase hypothesises that disease conditions occupying distinct anatomical regions can be factorised in the latent space under conditional assumptions over shared anatomy and disease-specific effects.
Under that setting, compositional generation should produce clinically plausible co-morbid samples even when the corresponding joint presentation is rare or absent in the training data.

This claim is evaluated along four complementary dimensions:
\begin{itemize}
    \item \textbf{Pathology presence:} dedicated pathology classifiers verify that both target disease conditions are present in the generated image, providing the clinical analogue of factor recovery in Phase~2.

    \item \textbf{Spatial plausibility:} pathological features appear in anatomically appropriate regions and do not interfere with one another, testing whether factorisation preserves non-conflicting disease evidence.

    \item \textbf{Radiological fidelity:} a classifier trained on real co-morbid cases assesses whether the generated image is consistent with genuine clinical presentations.

    \item \textbf{OOD setup:} co-morbid presentations are absent or severely underrepresented during training, ensuring that evaluation measures genuine out-of-distribution composition.
\end{itemize}

This phase matters because real co-morbid cases are rare and expensive to annotate in clinical datasets.
If successful, this phase would show that logical composition can generate clinically plausible examples of underrepresented combinations.
This would make it easier to study rare disease patterns and support data augmentation and downstream tasks in low-resource settings.




\begingroup
\captionsetup{skip=4pt}
\setlength{\floatsep}{6pt plus 1pt minus 2pt}
\setlength{\textfloatsep}{8pt plus 2pt minus 2pt}
\setlength{\intextsep}{8pt plus 2pt minus 2pt}

\subsection{Dataset Selection}
This phase requires a clinical dataset that is suitable for studying co-morbid composition under realistic conditions.
In particular, it must provide a sufficiently large set of posteroanterior (PA) chest X-rays, reliable pathology annotations, and enough variation in disease prevalence to reflect the low-resource nature of rare co-morbidity.
The VinBigData Chest X-ray dataset is well suited to this purpose because it contains 15{,}000 PA radiographs annotated by professional radiologists across 14 pathological findings.

As illustrated in Figure~\ref{fig:class_dist}, the data follows a pronounced long-tail distribution, with ``No finding'' accounting for 70.7\% of all images.
The remaining pathology classes are much less frequent, which creates a realistic setting in which disease-specific signals can easily be dominated by the normal findings.
For that reason, the training of pathology-specific experts will use balanced sampling so that rare disease patterns are not suppressed during learning.

\begin{figure}[ht]
\centering
\includegraphics[width=\textwidth]{chapters/research_method/media/xray/class_distribution.png}
\caption{Class distribution across 15,000 unique chest X-rays. The dominance of ``No finding'' (70.7\%) motivates the use of balanced sampling across pathology classes.}
\label{fig:class_dist}
\end{figure}

Within this dataset, we focus on cardiomegaly and pleural thickening as the primary pathology pair for Phase~3.
This pair is suitable because the two conditions do co-occur in clinical practice, yet their dominant visual manifestations are often concentrated in different anatomical regions.
Cardiomegaly primarily affects the enlarged cardiac silhouette in the central chest, whereas pleural thickening is expressed along the pleural boundary and lung periphery.
This makes them a useful test case for compositional generation, since their approximate spatial separation provides a plausible basis for the conditional assumptions needed to compose disease-specific factors without strong interference.

Their joint presentation, however, is much rarer in the training data than either pathology on its own.
As shown in Figure~\ref{fig:cooccurrence}, Cardiomegaly appears in 2,300 images and Pleural Thickening in 1,981, but the pair is annotated together in only 858 cases.
This gap is what makes the setting both clinically meaningful and methodologically challenging: a standard monolithic model sees many examples of each disease in isolation, but far fewer examples of how they appear together.
Our compositional approach instead asks whether the abundant single-disease cases can be used as separate priors and combined at inference time to synthesise plausible co-morbid presentations without requiring a large jointly labelled cohort.

\begin{figure}[H]
\centering
\includegraphics[width=0.63\textwidth]{chapters/research_method/media/xray/cooccurrence_heatmap_single.png}
\caption{VinBigData co-occurrence and scarcity heatmap. The bright diagonal reveals individual disease prevalence, while the dark off-diagonal cells highlight the low-resource co-morbidity problem.}
\label{fig:cooccurrence}
\end{figure}

\subsection{Experiment Design}

\subsubsection{Preprocessing}


Phase~3 begins by preparing the raw VinBigData data for modelling.
The dataset provides PA chest X-ray images in DICOM format together with annotation tables that specify pathology labels and bounding-box coordinates for localised findings.
We first match each radiograph to its corresponding annotation records using the image identifiers provided in the dataset metadata.

The raw images are stored as DICOM radiographs rather than standard image files.
This matters because DICOM preserves acquisition-dependent intensity conventions and scanner metadata that are not directly comparable across sites.
In particular, the \texttt{Photometric Interpretation} field determines whether brighter pixel values correspond to denser tissue or to air, so the same anatomy may appear with inverted intensity across images if this is not corrected.
To place all radiographs in a common intensity convention, we correct these inversions before further processing.

After inversion correction, the images are standardised to a common training representation.
If required by the model, each radiograph is resized to a fixed spatial resolution, and global pixel-intensity normalisation is applied to reduce residual contrast and brightness variation across scanners and clinical sites.
These steps reduce nuisance variation introduced during acquisition and make disease-related structure more comparable across the dataset.

Once the radiographs have been standardised, the bounding-box annotations become useful in two ways.
First, they provide an exploratory view of where the selected pathologies are typically expressed in the image.
As shown in Figure~\ref{fig:disease_samples}, cardiomegaly is concentrated around the enlarged cardiac silhouette in the central chest, while pleural thickening appears along the lung periphery and lower pleural margins.
Second, they provide spatial supervision for the model by identifying the regions where disease-specific changes should be represented while discouraging unnecessary changes elsewhere in the image.

\begin{figure}[ht]
\centering
\includegraphics[width=0.56\textwidth]{chapters/research_method/media/xray/disease_samples_bbox.png}
\caption{Ground-truth samples illustrating the spatial distinctness of Cardiomegaly (central) and Pleural Thickening (peripheral).}
\label{fig:disease_samples}
\end{figure}


\subsubsection{Conditional Assumptions}
\label{subsec:p3_assumptions}

We test the claim that out-of-distribution co-morbid chest X-rays can be generated when the learned representation separates shared anatomy from disease-specific effects.
This is the clinical extension of the idea developed in Phase~2: if improving the representation makes composition more reliable in a controlled setting, then a suitably factorised representation may also support composition in real medical images.
To evaluate that claim, we must state the conditions under which the score-addition formula can be used as a reasonable approximation in this setting.


Thus, the score-addition formula
\begin{equation}
    p(x \mid c_1, c_2)
    \;\approx\;
    \frac{p(x \mid c_1)\,p(x \mid c_2)}{p_0(x)}
    \label{eq:score_composition}
\end{equation}
is valid only when three assumptions hold jointly.
In Phase~3, these assumptions define the conditions under which compositional generation is expected to produce clinically plausible co-morbid chest X-rays rather than unrealistic superpositions.

\paragraph{Assumption 1: Separating Common Factors Of Variation From Disease Factors.}
Both single-disease distributions can be written as perturbations of a common underlying distribution:
\begin{equation}
    p(x \mid c_1) \approx p_0(x)\,\phi_1(x),
    \qquad
    p(x \mid c_2) \approx p_0(x)\,\phi_2(x),
\end{equation}
so that each disease-specific factor modifies the same underlying anatomy rather than redefining the entire image distribution.
In this setting, $p_0(x)$ captures chest anatomy, projection geometry, scanner characteristics, and non-pathological variation.
The ``No finding'' cases in VinBigData provide exactly this base distribution.
This motivates the use of a shared latent component $z_c$ to represent background anatomy and non-disease variation, rather than allowing disease-specific structure to absorb it.

\paragraph{Assumption 2: Disease Factors Do Not Interfere.}
Once shared anatomy is accounted for, the image evidence for $c_1$ and $c_2$ can be composed without major contradiction.
In latent terms,
\begin{equation}
    z = (z_c,\, z_{d1},\, z_{d2}),
    \qquad
    z_{d1} \perp z_{d2} \mid z_c,
    \label{eq:latent_factorisation}
\end{equation}
with $z_{d1}$ controlling only disease-1 structure and $z_{d2}$ controlling only disease-2 structure.
This is a stronger requirement than population-level independence of the diseases as labels: it requires that, conditional on shared anatomy, the image evidence for each disease adds without conflicting.

\paragraph{Assumption 3: Non-Overlapping Diseases Are Approximately Separable.}
The primary manifestations of each disease should be concentrated in regions that do not substantially overlap, \emph{and} the generative changes associated with one disease should not systematically alter the visual cues of the other.
Spatial separation supports this assumption but does not alone establish it: cardiomegaly acts mainly through heart-region structure, such as the cardiothoracic ratio and cardiac silhouette, while pleural thickening acts mainly through pleural boundary structure, such as peripheral pleural opacity and lower pleural margin changes.
These are distinct anatomical substrates with limited shared feature statistics, which is why spatial orthogonality provides practical evidence for, rather than a guarantee of, Assumption~2.

This assumption is motivated by the spatial distribution of the selected pathologies in VinBigData.
As illustrated in Figures~\ref{fig:density_maps} and~\ref{fig:spatial_comparison}, the two diseases can co-occur clinically, but their dominant visual manifestations are often concentrated in different parts of the chest.
The aggregate density maps indicate that cardiomegaly is concentrated in the central heart region, whereas pleural thickening is concentrated along the peripheral lower pleural margins.
The clinical example in Figure~\ref{fig:spatial_comparison} shows the same pattern at the level of an individual image, where the cardiac and pleural regions remain physically separated.

\begin{figure}[!tbp]
\centering
\begin{minipage}{0.48\textwidth}
\centering
\includegraphics[width=\linewidth]{chapters/research_method/media/xray/spatial_orthogonality.png}
\caption{Aggregate annotation density for Cardiomegaly ($n=2{,}300$) and Pleural Thickening ($n=1{,}981$), showing distinct hotspots in the central heart region and along the peripheral lower pleural margins respectively.}
\label{fig:density_maps}
\end{minipage}\hfill
\begin{minipage}{0.48\textwidth}
\centering
\includegraphics[width=\linewidth]{chapters/research_method/media/xray/normal_vs_comorbid.png}
\caption{Comparison between a normal lung and a co-morbid patient. The separation of the cardiac (red) and pleural (blue) bounding regions illustrates the anatomical distinctness of the two conditions.}
\label{fig:spatial_comparison}
\end{minipage}
\end{figure}

Together, Assumptions~1--3 reframe the H3 claim precisely: composition succeeds not because cardiomegaly and pleural thickening happen to occupy different parts of the image, but because there exists a representation in which shared anatomy is separated from disease-specific factors, disease effects are localised and approximately non-interfering, and composition in that representation remains on the plausible chest X-ray manifold.
\subsubsection{Architecture}

Phase~3 extends the partitioned representation strategy introduced in Regime~D of Phase~2 to the clinical setting.
The central idea is unchanged: composition is more likely to succeed if shared context is separated from disease-specific variation before inference-time combination is attempted.

We treat pathologies as disentangled ``objects'' and their visual features as ``attributes'' to formally model the spatial and functional relationships required for co-morbid representation.
By explicitly enforcing factorised representations, we aim to separate shared chest anatomy from disease-specific effects so that pathological features can be combined without collapsing into a single entangled signal.

The model therefore uses a partitioned latent code
\begin{equation}
    z = (z_c,\, z_{d1},\, z_{d2}),
\end{equation}
where $z_c$ captures shared anatomy, acquisition context, and other non-pathological variation, while $z_{d1}$ and $z_{d2}$ capture the two disease-specific factors.
This is the clinical analogue of the Regime~D design in Phase~2, with the latent partition now aligned to shared anatomy and the selected pathology pair rather than to synthetic attributes.

\subsubsection{Objective Functions}

The training objective is designed to make the latent partition meaningful rather than merely symbolic.
Let $y_1$ and $y_2$ denote the labels for the selected pathology pair, and let $M_1$ and $M_2$ denote the corresponding binary masks derived from the bounding-box annotations when present.
This objective combines four ingredients:
\begin{itemize}
    \item \textbf{Reconstruction and KL regularisation:} these preserve faithful reconstruction while keeping the shared and disease-specific latent blocks well behaved.

    \item \textbf{Predictiveness and exclusivity:} these terms encourage each disease-specific latent part to encode its intended pathology while discouraging the same signal from leaking into the opposite disease latent.

    \item \textbf{Dependence reduction:} mutual-information and orthogonality penalties reduce residual dependence and geometric overlap between $z_{d1}$ and $z_{d2}$.

    \item \textbf{Region-weighted reconstruction:} bounding-box-guided reconstruction promotes locality by forcing disease-related changes to be explained primarily within the annotated pathology regions.
\end{itemize}
thus the equation for the objective function is:
\begin{equation}
\begin{aligned}
\mathcal{L}_{\mathrm{P3}}
&=
\mathcal{L}_{\mathrm{rec}}^{\mathrm{rw}}(x,\hat{x};M_1,M_2)
+ \beta_c \, D_{\mathrm{KL}}\!\left(q_\phi(z_c \mid x)\,\|\,p(z_c)\right)
+ \beta_{d1} \, D_{\mathrm{KL}}\!\left(q_\phi(z_{d1} \mid x)\,\|\,p(z_{d1})\right) \\
&\quad
+ \beta_{d2} \, D_{\mathrm{KL}}\!\left(q_\phi(z_{d2} \mid x)\,\|\,p(z_{d2})\right)
+ \lambda^{+}_{1}\,\mathrm{CE}\!\left(g_{1}(\mu_{d1}), y_{1}\right)
+ \lambda^{+}_{2}\,\mathrm{CE}\!\left(g_{2}(\mu_{d2}), y_{2}\right) \\
&\quad
+ \lambda^{-}_{1}\,\mathrm{CE}\!\left(g_{1}(\mathrm{GRL}(\mu_{d2})), y_{1}\right)
+ \lambda^{-}_{2}\,\mathrm{CE}\!\left(g_{2}(\mathrm{GRL}(\mu_{d1})), y_{2}\right) \\
&\quad
+ \lambda_{\mathrm{mi}} \, \hat{I}(z_{d1}; z_{d2})
+ \lambda_{\mathrm{orth}} \, \left\lVert \mathrm{Cov}(z_{d1}, z_{d2}) \right\rVert_F^2 .
\end{aligned}
\end{equation}
Here, $\mathcal{L}_{\mathrm{rec}}^{\mathrm{rw}}$ is a region-weighted reconstruction term that places greater emphasis on the annotated pathology regions while discouraging unnecessary changes elsewhere in the image. A simple form is
\begin{equation}
\mathcal{L}_{\mathrm{rec}}^{\mathrm{rw}}(x,\hat{x};M_1,M_2)
=
\sum_i w_i \, \ell(x_i,\hat{x}_i),
\qquad
w_i =
\begin{cases}
w_{\mathrm{fg}}, & i \in M_1 \cup M_2, \\
w_{\mathrm{bg}}, & i \notin M_1 \cup M_2,
\end{cases}
\quad
w_{\mathrm{fg}} > w_{\mathrm{bg}},
\end{equation}
where $\ell(\cdot,\cdot)$ is the chosen pixelwise reconstruction loss. 

Together, these terms encourage a representation in which shared anatomy is preserved in $z_c$, disease-specific signal is localised to the appropriate latent part, and the two disease factors can later be composed without strong interference.


\subsection{Evaluation}
\label{subsec:p3_eval}

Phase~3 evaluates whether a factorised clinical representation supports held-out co-morbid composition under realistic conditions.
The evaluation mirrors the logic of Phase~2, but replaces synthetic attribute recovery with clinical criteria.
Each composed image is assessed along the following dimensions:
\begin{itemize}
    \item \textbf{Pathology presence:} two external pathology classifiers, trained on data splits disjoint from generative-model training, are used as the clinical analogue of factor-recovery tests.
    Given a composed image $x_{\mathrm{comp}}$, the cardiomegaly classifier $f_{\mathrm{card}}$ and the pleural-thickening classifier $f_{\mathrm{pleural}}$ each output a probability of presence.
    A composition is counted as pathology-preserving when both classifiers exceed predefined decision thresholds, indicating that the generated image contains recoverable evidence for both target diseases.

    \item \textbf{Co-morbidity fidelity:} a third classifier, $f_{\mathrm{co}}$, is trained on real co-morbid chest X-rays to distinguish genuine joint presentations from non-co-morbid cases.
    Its role is not to check whether both labels are merely present, but whether the generated image remains on the manifold of plausible co-morbid radiographs rather than appearing as an artificial overlay of isolated findings.
    High fidelity scores therefore provide a complementary measure of radiological realism beyond simple pathology presence.

    \item \textbf{Spatial plausibility:} spatial plausibility is assessed with the disease-specific anatomical regions derived from the VinBigData bounding boxes.
    The dominant abnormal signal in a composed image should align with the expected heart-region territory for cardiomegaly and the expected peripheral pleural territory for pleural thickening, without substantial cross-region leakage.
    These checks operationalise the requirement that the two disease factors remain anatomically distinct rather than collapsing into a spatially incoherent hybrid.

    \item \textbf{Held-out composition setting:} the final evaluation is performed in the low-resource regime introduced in Dataset Selection, where joint cardiomegaly--pleural-thickening cases are absent or severely underrepresented during generative training, while single-disease cases remain available.
    At inference time, disease-specific factors are composed in latent space and decoded into a candidate co-morbid radiograph.
    Composition succeeds when the generated sample is scored positive by both pathology classifiers, receives a high co-morbidity fidelity score, and satisfies the spatial plausibility check.
\end{itemize}

\subsection{Summary}

Phase~3 transfers the representation hypothesis of Phase~2 to a clinical domain in which factor independence is only approximate rather than engineered.
VinBigData provides a realistic low-resource setting: the selected pathologies are individually common enough to learn separately, yet their joint presentation remains relatively scarce.
The experiment design therefore standardises raw DICOM radiographs, uses annotation-derived spatial evidence to motivate the compositional assumptions, and extends the partitioned latent strategy of Regime~D so that shared anatomy is separated from disease-specific effects.

The central test is whether this representation supports held-out co-morbid composition.
If composed samples recover both target pathologies, remain spatially plausible, and are judged consistent with real co-morbid radiographs, then Phase~3 supports H3 by showing that approximate anatomical separability is sufficient for clinically plausible out-of-distribution composition.
If these criteria are not met, then Phase~3 does not support H3: either the conditional assumptions are too weak in this clinical setting, or the learned representation is still not sufficiently factorised to support stable composition.